\section{First Duplicate: a Search with an Early Exit}\label{sec:fd}

\subsection{The program}
\begin{lstlisting}
method main(a: array<int, int>, n: int) returns (r: bool)
  requires 0 <= n
  ensures r ==> !distincts(selects(a, butlast(counter(0, 1, natural(n)))))
  ensures !r ==> distincts(selects(a, butlast(counter(0, 1, natural(n)))))
{
  var i: int; var j: int;
  r := false; i := 0;
  while i < n && !r {
    j := i + 1;
    while j < n && !r {
      if a[i] == a[j] { r := true; }
      j := j + 1;
    }
    i := i + 1;
  }
}
\end{lstlisting}
The nested loop of \texttt{distinct.dfy} (which is verified) with both loops stopped by the flag: the
search for a duplicate ends at the first one. The postcondition uses $\mathsf{distincts}$ of the lemma
database, and is stated as two implications rather than the equivalence $r \Leftrightarrow \neg\mathsf{distincts}(\ldots)$,
for a reason explained below.

\subsection{The trace equation system}
The outer loop's constraint includes the flag, $\al{\bl{K}, \bl{R}, \bl{I}}{k, r, i \mapsto \neg r}$, so
$r$ is false at the head of every outer iteration, and the simplifier uses this inside the inner trace:
\begin{align}
  K &= I = \cnt(0, 1, n') \qquad r = \cur(R) \qquad R = \mathsf{results}(\mathsf{false}, \bl{K}, \seq{R}) \qquad i = \intf(n') \label{eq:fd0}\\
  &\al{\bl{K}, \bl{R}, \bl{I}}{k, r, i \mapsto \neg r} \label{eq:fd1}\\
  &\al{\bl{K}, \bl{I}, \seq{R}, \seq{R}'}{k, i, r, R' \mapsto r = \cur(R')} \label{eq:fd2}\\
  &\al{\bl{K}, \bl{I}, \seq{R}', \seq{J}'}{k, i, R', J \mapsto {}\nonumber\\
  &\qquad\qquad R' = \tr{\mathsf{false}}{\bl{J}}{r, j \mapsto \upd{a[k] = a[j]}{\mathsf{true}}\ \ \upd{a[k] \ne a[j]}{\mathsf{false}}}} \label{eq:fd3}\\
  &\al{\bl{K}, \seq{N}, \seq{J}'}{k, n'', J \mapsto J = \cnt(k + 1, 1, n'')} \label{eq:fd4}\\
  &\al{\bl{K}, \seq{N}, \seq{J}}{k, n'', j \mapsto j = \intf(n'') + k + 1} \label{eq:fd4b}\\
  &\al{\bl{K}, \bl{I}, \seq{J}'}{k, i, J \mapsto \al{\bl{J}}{j \mapsto j < n}} \label{eq:fd5}\\
  &\al{\bl{K}, \bl{I}, \seq{R}'}{k, i, R' \mapsto \al{\bl{R'}}{r \mapsto \neg r}} \label{eq:fd5b}\\
  &\al{\bl{K}, \bl{I}, \seq{R}, \seq{J}}{k, i, r, j \mapsto \neg(j < n \wedge \neg r)} \label{eq:fd6}\\
  &\neg(i < n \wedge \neg r)\ \Longrightarrow\ r \Leftrightarrow \neg\mathsf{distincts}(\mathsf{selects}(a, \bl{\cnt(0, 1, \natf(n))})) \nonumber
\end{align}
(after the fragment lemmas of Section~\ref{sec:sel:frag}, which gave \eqref{eq:fd4} and \eqref{eq:fd4b}). The inner update
\eqref{eq:fd3} does not mention the previous value of the flag, since the inner constraint
$\neg r$ \eqref{eq:fd5b} replaced it by $\mathsf{false}$: the inner trace is a sequence of comparisons
with $a[k]$. The inner exit \eqref{eq:fd6} is the negation of the test $j < n \wedge \neg r$: either all
later positions were visited, or the flag was set. The outer exit $\neg(i < n \wedge \neg r)$ is a
disjunction, $n \le i \vee r$, with one goal per postcondition.

\subsection{The lemmas}
The summary of the inner loop in the style of Section~\ref{sec:sel:sum} is a rewrite of the fragment \eqref{eq:fd2}
into a fact about the results: with \eqref{eq:fd3}--\eqref{eq:fd6} as
side conditions (duplicated and identified by equations where closed),
\begin{align*}
  \al{ks, rs, Rs}{k, r, R' \mapsto r = \cur(R')}\ \Longleftrightarrow\ \al{ks, rs}{k, r \mapsto r = \neg\mathsf{nodup}(a, k, n)}
\end{align*}
where $\mathsf{nodup}(a, k, n)$ is the abstract "$a[k]$ does not occur among $a[k + 1..n)$". This is
sound in both exit cases of the inner loop: if the flag was set, a later duplicate was found; if $j$
reached $n$, all later positions were compared. The outer trace is then
$\mathsf{results}(\mathsf{false}, \bl{K}, \seq{R})$ with $\seq{R}$ characterized pointwise, and the
simplifier itself turns it into $\mathsf{results}(\mathsf{false}, \bl{K}, \mapc{\bl{K}}{k \mapsto \neg\mathsf{nodup}(a, k, n)})$.

The outer loop is summarized on the atom $\mathsf{distincts}(\ldots)$ of the postcondition, by one rule per
case of the exit. If the flag is set, it is the result of the last iteration, which found a duplicate:
\begin{align*}
  &\lemma{distincts\_flag\_set}:\quad Kv = \cnt(0, 1, m)\\
  &\quad\wedge\ flag = \cur(\mathsf{results}(\mathsf{false}, \bl{Kv}, \mapc{\bl{Kv}}{k \mapsto \neg\mathsf{nodup}(a, k, n)}))\\
  &\quad\wedge\ flag\ \wedge\ \intf(m) \le n\ \Longrightarrow\ \mathsf{distincts}(\mathsf{selects}(a, \bl{\cnt(0, 1, \natf(n))}))\ \Longleftrightarrow\ \mathsf{false}
\end{align*}
If all iterations ran, $n \le \intf(m)$, and the flag was reset at the head of each, \eqref{eq:fd1}, then no
element before the last has a later duplicate (the flag at the head of iteration $k$ is the result of
iteration $k - 1$) and the last has none:
\begin{align*}
  &\lemma{distincts\_flags\_reset}:\quad Kv = \cnt(0, 1, m)\ \wedge\ Rv = \mathsf{results}(\mathsf{false}, \bl{Kv}, rs)\\
  &\quad\wedge\ \al{\bl{Kv}, hs}{k, h \mapsto \neg h}\ \wedge\ hs = \bl{Rv}\\
  &\quad\wedge\ \al{\bl{Kv}, rs_2}{k, h \mapsto h = \neg\mathsf{nodup}(a, k, n)}\ \wedge\ rs_2 = rs\ \wedge\ n \le \intf(m)\\
  &\quad\Longrightarrow\ \mathsf{distincts}(\mathsf{selects}(a, \bl{\cnt(0, 1, \natf(n))}))\ \Longleftrightarrow\ \mathsf{true}
\end{align*}
Three details made these fire. The outer counter and the flag trace are the variables $K$ and $R$ of the
goal with their equations \eqref{eq:fd0} as hypotheses, so a rule with $\bl{\cnt(0, 1, m)}$ as an
\emph{argument} never matches; the open trace variables $Kv$, $Rv$ are matched against the hypotheses
(Section~\ref{sec:merge}). The fact of the flags at the heads \eqref{eq:fd1} carries the counter $\bl{I}$ as
an unused dependency, which the simplifier drops from a closed instance of the condition, after which it
is not found among the facts; hence the open list $hs$, identified with $\bl{Rv}$ afterwards. And the exit
$n \le i \vee r$ is needed as a closed condition: with the postcondition as an equivalence, the goal's
conclusion is $i < n \wedge \neg r\ \vee\ r \Leftrightarrow \neg\mathsf{distincts}(\ldots)$, which
\texttt{-split} does not split, and in the context of the second disjunct neither literal of the negated
first is available, so the condition is not decided; probes (rules with the real left-hand side and the
first $k$ conditions) located it as the only failing one. As two implications, each goal has its literal
of the exit, $r$ respectively $n \le i$ with $i = \intf(n')$, in the context of the atom.

An earlier set of lemmas closed both goals and also the negated postcondition: the rule for the first step
had the summary fragment as an open side condition, and a rule $\mathsf{nodup}(a, k, n) \Leftrightarrow \mathsf{true}$
under the hypothesis $\mathsf{distincts}(\ldots)$ for the second. An $\mathsf{all}$ side condition is matched
in the simplifier's context of the fact's combinator, where rules fire: under the hypothesis $\mathsf{distincts}$,
the rule rewrote the condition to \eqref{eq:fd1}, a fact, so $\mathsf{distincts}$ was derived from itself
(finding 13 of \texttt{NOTES.md}). The rules above take the atom $\mathsf{distincts}$ as left-hand side
and never rewrite $\mathsf{nodup}$.

\subsection{The idea}
An early exit is a data-dependent number of iterations with a disjunctive exit; the summary of the
inner loop must be stated so that it holds in both cases, which the abstract atom $\mathsf{nodup}$
affords, since it describes the result rather than the number of iterations; the outer loop is summarized
on the atom of the postcondition with one rule per case of its exit, each with the case as a closed
condition. The soundness lesson stands: a side condition that other rules can simplify under the conclusion
to be proved is a circular justification, and the safe form of side conditions is an atom that is matched
literally (as \lemma{consts\_prefix\_def} of the database does for a constraint, and as the summary itself
does for the inner loop).
